% !TeX root = ../main.tex
\chapter{KẾT LUẬN VÀ HƯỚNG PHÁT TRIỂN}
\label{chap:ketluan}

\section{Tổng kết}

Luận văn đã xây dựng và đánh giá một quy trình học giám sát yếu hoàn chỉnh cho bài
toán phát hiện và phân đoạn bệnh lá sầu riêng, đi từ phân loại ảnh qua giải thích
bằng bản đồ kích hoạt, sinh nhãn giả, tinh chỉnh biên bằng mô hình nền tảng, đến
huấn luyện mô hình phân đoạn --- toàn bộ mà không sử dụng bất kỳ nhãn thủ công nào
ở mức điểm ảnh.

Kết quả chính được tóm tắt ở \cref{tab:tongket}.

\begin{table}[H]
\centering
\caption{Tổng hợp kết quả chính của luận văn}
\label{tab:tongket}
\small
\begin{tabular}{L{5.2cm} L{8.0cm}}
\toprule
\textbf{Hạng mục} & \textbf{Kết quả} \\
\midrule
Phân loại năm lớp bệnh & Accuracy 97,19\%, $F_1$ 97,19\% trên 890 ảnh kiểm thử \\
Độ chính xác kiểm định tốt nhất & 98,19\% tại epoch 11 \\
Sinh nhãn giả & 3.104 mặt nạ, không dùng nhãn thủ công nào \\
Kiểm soát độ phủ (V2) & Sai lệch dưới 0,4 điểm phần trăm mọi lớp \\
Loại nhãn nhiễu lớp lá khoẻ & 683 mặt nạ dương tính giả bị loại bỏ \\
Hiệu quả tinh chỉnh bằng SAM & IoU tăng 21,5\% (cận trên), hội tụ sớm hơn 8 epoch \\
Phát hiện về cơ chế bảo vệ IoU & Chứng minh không thể phát hiện hiện tượng nở mặt nạ \\
Phát hiện về tiêu chí dừng sớm & Mất 0,048--0,050 IoU khi dùng FocalDice \\
\bottomrule
\end{tabular}
\end{table}

\section{Trả lời các câu hỏi nghiên cứu}

\subsection{CH1 --- Hiệu năng phân loại và các cặp lớp dễ nhầm}

\emph{Câu hỏi: Mô hình \EffNet{} đạt độ chính xác bao nhiêu trên năm lớp bệnh lá
sầu riêng, và những cặp lớp nào dễ bị nhầm lẫn nhất?}

Mô hình đạt \textbf{97,19\% độ chính xác} trên tập kiểm thử 890 ảnh, với $F_1$ có
trọng số 97,19\%. Độ chính xác kiểm định tốt nhất là 98,19\% tại epoch 11; khoảng
cách 1,00 điểm phần trăm giữa hai giá trị cho thấy mô hình khái quát tốt. Cả năm
lớp đều đạt $F_1$ trên 95\%, biên độ giữa lớp cao nhất và thấp nhất chỉ 2,66 điểm
phần trăm.

Cặp nhầm lẫn chính là \textbf{Algal Leaf Spot và Phomopsis Leaf Spot}, chiếm 5
trên tổng số 24 lỗi. Kết quả này đã được dự báo từ giai đoạn phân tích khám phá dữ
liệu, dựa trên hai bằng chứng độc lập: quan sát hình thái cho thấy cả hai bệnh đều
tạo đốm nhỏ trên nền lá xanh, và phân tích màu cho thấy hai lớp có phân bố kênh
màu chồng lấn nhiều nhất.

Đáng chú ý về mặt ứng dụng, tỉ lệ bỏ sót bệnh --- ảnh bệnh bị phân loại thành lá
khoẻ --- chỉ là 6 trên 694 ảnh, tức 0,86\%.

\subsection{CH2 --- Chất lượng của bản đồ chú ý}

\emph{Câu hỏi: Bản đồ chú ý sinh ra từ mô hình phân loại có tập trung vào vùng tổn
thương thực sự hay không, và mức độ tin cậy của chúng khác nhau thế nào giữa các
lớp bệnh?}

Câu trả lời phân hoá rõ rệt theo lớp, và chính sự phân hoá này là kết quả có giá
trị nhất.

Với \textbf{bốn lớp bệnh}, bản đồ chú ý tập trung đúng vào vùng tổn thương: kích
hoạt bám vào các đốm với Algal và Phomopsis, bám vào vết cắn với Allocaridara, và
trải theo mảng hoại tử với Leaf Blight. Độ tin cậy trung bình trên tập huấn luyện
đạt 0,971--0,991.

Với \textbf{lớp lá khoẻ}, bản đồ chú ý \emph{khuếch tán khắp phiến lá}, không có
vùng trọng tâm. Hệ quả định lượng là mặt nạ sinh ra từ nhãn giả V1 cho lớp này có
độ phủ trung bình 26,4\% --- cao gần gấp đôi lớp bệnh lan rộng nhất --- với trường
hợp cực đoan lên tới 59,1\% trên một chiếc lá hoàn toàn khoẻ mạnh.

Đây là bằng chứng thực nghiệm trực tiếp cho luận điểm lý thuyết trung tâm của luận
văn: \textbf{bản đồ chú ý trả lời câu hỏi ``mô hình nhìn vào đâu để phân loại'',
không phải câu hỏi ``vùng bệnh nằm ở đâu''}. Với lá khoẻ, mô hình nhận diện qua sự
\emph{vắng mặt} của dấu hiệu bệnh --- một hành vi hoàn toàn hợp lý cho nhiệm vụ
phân loại nhưng vô nghĩa khi chuyển thành mặt nạ tổn thương.

Nhận thức này dẫn tới biện pháp ép mặt nạ lớp lá khoẻ về không, loại bỏ 683 mặt nạ
nhiễu và cung cấp cho mô hình phân đoạn các mẫu âm sạch.

\subsection{CH3 --- Hiệu quả của việc tinh chỉnh biên bằng SAM}

\emph{Câu hỏi: Việc tinh chỉnh biên bằng \SAM{} có cải thiện chất lượng nhãn giả,
và qua đó cải thiện mô hình phân đoạn, hay không?}

Câu trả lời có hai vế đối lập nhau và cả hai đều đúng.

\textbf{Xét theo tiêu chí bệnh học, \SAM{} làm nhãn xấu đi.} Thay vì chỉ tinh chỉnh
biên như kỳ vọng, \SAM{} mở rộng độ phủ mặt nạ trung bình từ 15,04\% lên 25,43\%
--- tăng 69,1\%, với mức tăng theo lớp từ 41,5\% đến 95,7\%. Có 200 trên 3.104 mặt
nạ vượt quá 50\% diện tích ảnh, điều bất khả thi với bệnh đốm lá. Nguyên nhân là
\SAM{} được huấn luyện để phân đoạn đơn vị ngữ nghĩa tự nhiên, mà đơn vị đó với
gợi ý điểm đặt trên đốm bệnh thường là \emph{cả chiếc lá}.

\textbf{Xét theo hiệu quả huấn luyện, \SAM{} cải thiện rõ rệt.} Với kiến trúc, hàm
mất mát và tập kiểm thử giống hệt nhau, phiên bản V3 vượt V2 trên toàn bộ bốn chỉ
số. Bằng chứng thuyết phục nhất không phải con số IoU mà là tốc độ hội tụ: V3 đạt
val\_loss 0,4478 tại epoch 12, trong khi V2 phải tới epoch 20 mới đạt 0,5957.

Hai vế này không mâu thuẫn vì chúng đo hai thứ khác nhau. Nhãn \SAM{} \emph{sai về
phạm vi} nhưng \emph{đúng về cấu trúc} --- biên bám theo ranh giới thực trong ảnh
và nhất quán giữa các ảnh. Mô hình học được từ tính nhất quán cấu trúc đó.

Về mức độ cải thiện, luận văn không kết luận toàn bộ 21,5\% là do \SAM{}. Ba yếu
tố gây nhiễu --- V2 bị dừng sớm khi còn đang cải thiện, thiên lệch của IoU theo
kích thước mục tiêu, và khác biệt kích thước lô --- đều tác động cùng chiều làm
phóng đại lợi thế của V3. Con số 21,5\% vì vậy là \textbf{cận trên} chứ không phải
giá trị thực.

\subsection{CH4 --- Khả năng so sánh của các chỉ số khi thiếu nhãn chuẩn}

\emph{Câu hỏi: Trong điều kiện không có nhãn chuẩn ở mức điểm ảnh, các chỉ số IoU
và Dice giữa những phiên bản khác nhau có so sánh trực tiếp được với nhau không?
Nếu không thì phép so sánh nào là hợp lệ?}

\textbf{Câu trả lời là không}, và đây là kết luận phương pháp luận quan trọng nhất
của luận văn.

Mỗi phiên bản được đánh giá trên mốc tham chiếu của chính nó: V1 so với nhãn V1,
V2 so với nhãn V2, V3 so với nhãn đã qua tinh chỉnh. Ba mốc tham chiếu khác nhau về
hình dạng, độ phủ lẫn độ biến thiên. So sánh trực tiếp các con số IoU giữa chúng
tương đương với việc so điểm của ba học sinh làm ba đề thi có độ khó khác nhau.

Minh hoạ rõ nhất là phiên bản V1 với IoU 0,7671 --- cao nhất trong ba phiên bản.
Con số này \emph{không} có nghĩa V1 định vị bệnh tốt nhất. Nó phản ánh việc mục
tiêu của V1 dễ hơn: nhãn V1 là các khối tròn lớn và trơn, hệ quả của việc phóng
bản đồ $7 \times 7$ lên $224 \times 224$, nên dễ tái tạo hơn nhiều so với các biên
sắc và phức tạp của nhãn \SAM{}.

Ngoài khác biệt về mốc tham chiếu, luận văn còn xác định thêm hai lý do kỹ thuật
khiến V1 không thể so sánh với V2 và V3: hai nhóm dùng \emph{hai tập kiểm thử khác
nhau} do khác bộ sinh số ngẫu nhiên dù cùng hạt giống, và cài đặt của V1 vô tình
áp dụng tăng cường dữ liệu lên cả tập kiểm định và kiểm thử. Tổng cộng cặp V1--V2
khác nhau ở \textbf{tám biến} cùng lúc.

\textbf{Phép so sánh hợp lệ duy nhất là V2 với V3}, vì hai phiên bản này dùng chung
kiến trúc, hàm mất mát, bộ tối ưu, bộ tăng cường và cùng một tập kiểm thử. Ngay cả
với cặp này, ba yếu tố gây nhiễu còn lại vẫn cần được trừ hao khi diễn giải.

Kết luận tổng quát: \textbf{trong khung học giám sát yếu không có nhãn chuẩn, chỉ
số IoU và Dice chỉ đo được mức độ mô hình tái tạo lại bộ nhãn giả của chính nó,
không đo được độ chính xác định vị bệnh tuyệt đối.} Mọi so sánh giữa các phiên bản
dùng bộ nhãn khác nhau đều cần được đặt trong ngữ cảnh này.

\section{Đóng góp của luận văn}

\subsection{Đóng góp về phương pháp}

\textbf{Cơ chế ngưỡng phân vị theo lớp.} Thay vì cố định giá trị ngưỡng và để độ
phủ trở thành kết quả không kiểm soát được, luận văn cố định tỉ lệ độ phủ và suy
ngược ra ngưỡng cho từng ảnh. Cơ chế này đạt độ phủ mục tiêu với sai lệch dưới 0,4
điểm phần trăm ở mọi lớp, đồng thời tránh được hiện tượng phân đoạn thừa của ngưỡng
Otsu khi bản đồ nhiệt hợp nhất đa tỉ lệ trở nên trơn và đơn đỉnh.

\textbf{Đưa tri thức tiên nghiệm vào quy trình sinh nhãn.} Việc ép mặt nạ lớp lá
khoẻ về không là một can thiệp đơn giản nhưng hiệu quả, loại bỏ 683 mặt nạ dương
tính giả mà bản thân bản đồ kích hoạt không bao giờ tự loại được.

\subsection{Đóng góp về phát hiện thực nghiệm}

\textbf{Chứng minh cơ chế bảo vệ dựa trên IoU không thể phát hiện hiện tượng nở
mặt nạ.} Khi mặt nạ mới chứa trọn mặt nạ gốc, IoU giữa chúng bằng đúng tỉ số hai
độ phủ. Do đó cơ chế bảo vệ với ngưỡng $\theta$ chỉ kích hoạt khi độ phủ mới vượt
$1/\theta$ lần độ phủ gốc. Với $\theta = 0{,}15$ và độ phủ thực tế 17,6--21,7\%,
ngưỡng cần thiết là 117--145\% diện tích ảnh --- bất khả thi về mặt vật lý. Đây là
khiếm khuyết thiết kế mang tính hệ thống, không phải lỗi cài đặt, và bắt nguồn từ
tính đối xứng của IoU: nó không phân biệt được ``lệch vị trí'' với ``đúng vị trí
nhưng quá rộng''. Luận văn đề xuất bổ sung một ràng buộc bất đối xứng về tỉ lệ độ
phủ để khắc phục.

\textbf{Xác định hệ quả của việc chọn điểm kiểm tra theo hàm mất mát khi dùng
FocalDice.} Với phiên bản V1 dùng BCE + Dice, tổn thất do chọn theo hàm mất mát
thay vì theo IoU chỉ là 0,0016 --- không đáng kể. Với hai phiên bản dùng FocalDice,
tổn thất lớn gấp ba mươi lần: 0,0478 và 0,0502. Nguyên nhân là thành phần Focal
thu nhỏ đóng góp của các điểm ảnh dễ, khiến giá trị hàm mất mát bị chi phối bởi
một nhóm nhỏ điểm ảnh biên trong khi IoU tính trên toàn vùng. Khuyến nghị rút ra
có thể áp dụng trực tiếp cho các nghiên cứu tương tự.

\subsection{Đóng góp về phương pháp luận}

\textbf{Khung phân tích tách bạch cho học giám sát yếu.} Luận văn đề xuất và áp
dụng cách tiếp cận đếm số biến khác nhau giữa các cặp cấu hình trước khi diễn giải
chênh lệch hiệu năng. Cách làm này cho thấy cặp V1--V2 khác nhau ở tám biến nên
không thể diễn giải, còn cặp V2--V3 khác ba biến nên diễn giải được với điều kiện
trừ hao. Đồng thời luận văn cũng đánh giá \emph{chiều tác động} của từng yếu tố
gây nhiễu để xác định con số quan sát được là cận trên hay cận dưới.

\textbf{Tài liệu kiểm kê số liệu.} Toàn bộ con số trong luận văn được truy vết về
đúng notebook và ô lệnh sinh ra nó, kèm danh sách các điểm chưa đủ bằng chứng.
Cách làm này bảo đảm tính minh bạch và cho phép kiểm chứng độc lập.

\section{Hạn chế}

\subsection{Hạn chế cốt lõi --- thiếu nhãn chuẩn}

Hạn chế bao trùm và quyết định phạm vi mọi kết luận là bộ dữ liệu không có nhãn
thủ công ở mức điểm ảnh. Hệ quả là:

\begin{itemize}
    \item Không thể xác lập hiệu năng định vị bệnh \emph{tuyệt đối} của bất kỳ
          phiên bản nào.
    \item Không thể hiệu chỉnh các tỉ lệ độ phủ mục tiêu 18--22\%, vốn hiện được
          chọn theo suy luận bệnh học.
    \item Không thể kiểm chứng liệu hiện tượng \SAM{} mở rộng mặt nạ có thực sự
          làm nhãn xa rời tổn thương thật hay không --- kết luận hiện tại dựa trên
          quan sát định tính và lập luận về diện tích hợp lý.
\end{itemize}

\subsection{Hạn chế về thiết kế thực nghiệm}

Các phiên bản không được thiết kế như một thí nghiệm có kiểm soát ngay từ đầu.
Cặp V1--V2 khác nhau ở tám biến; ngay cả cặp V2--V3 vẫn còn hai khác biệt ngoài ý
muốn về kích thước lô và patience. Ràng buộc bộ nhớ đồ hoạ 2,15 GB là nguyên nhân
trực tiếp của một số khác biệt này.

Ngoài ra, mỗi cấu hình chỉ được chạy \emph{một lần} với một hạt giống ngẫu nhiên
duy nhất. Không có ước lượng phương sai, nên không thể khẳng định chênh lệch quan
sát được vượt ngưỡng dao động ngẫu nhiên giữa các lần chạy.

\subsection{Hạn chế về dữ liệu}

Toàn bộ ảnh đã được co về $224 \times 224$ và bị xoá siêu dữ liệu EXIF trước khi
đến với nghiên cứu, nên không thể khảo sát ảnh hưởng của độ phân giải đầu vào hay
điều kiện chụp. Hệ thống cũng chỉ được đánh giá trên ảnh chụp lá đơn lẻ, chưa kiểm
chứng trên ảnh chụp toàn tán trong điều kiện đồng ruộng --- nơi có nền phức tạp,
nhiều lá chồng lấn và điều kiện ánh sáng biến thiên.

\subsection{Hạn chế về phạm vi bài toán}

Bài toán phân đoạn được đặt ở dạng nhị phân --- bệnh so với nền --- nên mô hình
phân đoạn không phân biệt được loại bệnh ở mức điểm ảnh. Với ảnh có nhiều loại
bệnh cùng lúc, hệ thống hiện chỉ cho một mặt nạ chung.

\section{Hướng phát triển}

\subsection{Ưu tiên cao}

\textbf{Xây dựng tập kiểm thử có nhãn chuẩn.} Đây là việc cần làm trước tiên vì nó
gỡ bỏ hạn chế cốt lõi. Chỉ cần một tập nhỏ 50--100 ảnh được chuyên gia khoanh vùng
thủ công là đủ để: đo hiệu năng định vị bệnh tuyệt đối, so sánh công bằng ba phiên
bản trên cùng một mốc tham chiếu, và hiệu chỉnh các tỉ lệ độ phủ mục tiêu. Chi phí
cho quy mô này hoàn toàn khả thi.

\textbf{Sửa cơ chế bảo vệ và chạy lại giai đoạn tinh chỉnh.} Bổ sung ràng buộc bất
đối xứng về tỉ lệ độ phủ như đã đề xuất, với hệ số $\kappa = 2{,}0$, rồi sinh lại
nhãn và huấn luyện lại. Đây là thí nghiệm rẻ nhưng cho phép kiểm chứng trực tiếp
liệu việc ngăn hiện tượng nở có cải thiện chất lượng nhãn hay không.

\textbf{Chuyển tiêu chí chọn điểm kiểm tra sang val\_IoU.} Chỉ là thay đổi vài
dòng mã nhưng theo số liệu đã đo sẽ thu hồi được khoảng 0,05 IoU cho phiên bản V3.

\subsection{Ưu tiên trung bình}

\textbf{Thiết kế lại thí nghiệm có kiểm soát chặt.} Cố định toàn bộ siêu tham số
giữa các phiên bản và chỉ thay đúng một biến mỗi lần, đồng thời chạy nhiều hạt
giống ngẫu nhiên để có ước lượng phương sai.

\textbf{So sánh \SAM{} với Dense CRF.} Dense CRF là hướng tinh chỉnh nhãn thay thế,
không cần mô hình nền tảng nặng và bám theo biên màu sắc thực. Hướng này đã được
thử nhưng không triển khai được do thư viện không tương thích với phiên bản Python
sử dụng; việc so sánh vẫn là câu hỏi bỏ ngỏ.

\textbf{Khảo sát cường độ tăng cường dữ liệu.} Kiểm chứng giả thuyết rằng các phép
biến dạng đàn hồi làm nhiễu tín hiệu huấn luyện khi nhãn vốn đã không chắc chắn về
biên.

\textbf{Thử \SAM{}2.} Phiên bản kế tiếp của \SAM{} có chất lượng biên tốt hơn và
có thể giảm mức độ nở mặt nạ.

\subsection{Hướng mở rộng dài hạn}

\textbf{Phân đoạn đa lớp.} Mở rộng từ nhị phân sang phân đoạn theo từng loại bệnh,
cho phép xử lý ảnh có nhiều bệnh đồng thời.

\textbf{Ước lượng mức độ bệnh.} Khắc phục hạn chế đã nêu rằng nhãn giả V2 không mã
hoá được mức độ nặng nhẹ, hướng tới đầu ra là tỉ lệ diện tích lá bị bệnh có hiệu
chỉnh --- đại lượng trực tiếp phục vụ quyết định canh tác.

\textbf{Kiểm chứng trong điều kiện đồng ruộng.} Đánh giá trên ảnh chụp toàn tán với
nền phức tạp và điều kiện ánh sáng biến thiên, đối chiếu với bài học từ bộ dữ liệu
PlantVillage về khoảng cách giữa dữ liệu chuẩn hoá và dữ liệu thực
tế~\cite{mohanty2016plantdisease}.

\textbf{Triển khai trên thiết bị biên.} Lượng tử hoá và cắt tỉa mô hình để chạy
trực tiếp trên điện thoại, phục vụ chẩn đoán tại vườn không cần kết nối mạng.

\section{Lời kết}

Luận văn cho thấy một quy trình học giám sát yếu hoàn toàn có thể chuyển nhãn mức
ảnh --- vốn rẻ và sẵn có --- thành mặt nạ phân đoạn ở mức điểm ảnh mà không cần bất
kỳ công khoanh vùng thủ công nào. Mô hình phân loại đạt 97,19\% độ chính xác, và
toàn bộ 3.104 mặt nạ giả được sinh tự động.

Nhưng đóng góp mà luận văn coi trọng nhất không nằm ở các con số hiệu năng. Đó là
việc chỉ ra rằng \emph{trong khung học giám sát yếu, chính các con số ấy cần được
đọc một cách thận trọng}. Phiên bản có IoU cao nhất không phải phiên bản định vị
bệnh tốt nhất. Một mô hình nền tảng có thể vừa làm hỏng nhãn theo tiêu chí bệnh học
vừa cải thiện kết quả huấn luyện. Và một cơ chế bảo vệ tưởng như hợp lý lại có thể
bất lực về mặt toán học trước đúng vấn đề mà nó được thiết kế để ngăn chặn.

Những phát hiện này chỉ lộ ra khi đối chiếu kỹ số liệu thực tế với kỳ vọng thiết
kế, thay vì dừng lại ở con số tốt nhất đạt được. Đó cũng là tinh thần mà luận văn
mong muốn đóng góp cho các nghiên cứu học giám sát yếu tiếp theo trong lĩnh vực
nông nghiệp.
